\documentclass[12pt,a4paper]{article}
\usepackage[utf8]{inputenc}
\usepackage[portuguese]{babel}
\usepackage[T1]{fontenc}
\usepackage{amsmath}
\usepackage{amsfonts}
\usepackage{amssymb}
\usepackage{graphicx}
\usepackage{booktabs}
\usepackage{hyperref}
\usepackage{geometry}
\geometry{margin=2.5cm}

\title{\textbf{Relatório sobre consumo de APIs}}
\author{}
\date{}

\begin{document}

\maketitle

\section{IBGE}

Essa API basicamente serve para extrair dados categóricos e numéricos sobre varias coisas relacionadas ao Brasil. Nessa API é possível pesquisar coisas como

\subsection{Áreas Urbanizadas}
12747- Áreas urbanizadas densas
12748- Áreas urbanizadas pouco densas
12749- Total de áreas urbanizadas
12750- Loteamento vazio
12751- Área total mapeada
12752- Outros equipamentos urbanos
12753- Vazios intraurbanos
12754- Vazios remanescentes


\end{document}